\documentclass[../../main.tex]{subfiles}
\begin{document}
\section{伪 Riemann 度量}\label{section:ch13-pseudo-riemannian}

通过放宽正定性要求，可得到 Riemann 度量的一个重要推广。向量空间 $V$ 上的一个对称 $2$-张量 $g$ 称为\emph{非退化的}，若由 $\hat g(v)(w)=g(v,w)$ 定义的线性映射 $\hat g:V\to V^*$ 是一个同构；等价地，若对每个非零 $v\in V$，都存在 $w\in V$ 使 $g(v,w)\ne0$。正如任何内积都可以通过换到一组标准正交基而化为 Euclid 内积一样，每个非退化对称 $2$-张量都可以通过基变换化为矩阵对角元全为 $\pm1$ 的形式（其证明是 Gram--Schmidt 算法的一个变形）。正、负对角元的个数 $r,s$ 与基的选取无关（这一事实称为 Sylvester 惯性律；证明见\cite{FIS03}）。于是有序对 $(r,s)$ 称为 $g$ 的\emph{符号差}，是 $g$ 的一个不变量。

光滑流形 $M$ 上的一个\emph{伪 Riemann 度量}是一个光滑对称 $2$-张量场，它在每一点处非退化，且在 $M$ 上各处符号差相同。符号差为 $(n-1,1)$（或 $(1,n-1)$，视所用约定而定）的伪 Riemann 度量称为\emph{Lorentz 度量}；它们在物理中起着核心作用，被用于 Einstein 广义相对论中对引力的建模。

我们不再进一步探讨伪 Riemann 度量这一主题，仅指出：Riemann 度量存在性的证明不能照搬到伪 Riemann 情形，因为带正系数的非退化 $2$-张量的线性组合一般未必非退化。事实上，并非每个流形都容许 Lorentz 度量（参见\cite{HE73}，第~39 页）。
\end{document}
